\newpage
\Article{69}
Moduli $a+bi$ for which $b=0$ does not hold (numbers of the fourth species in Art.~34) are much harder to deal with, however, and several investigations must come first. We shall denote the norm $aa+bb$, which will be a real prime of the form $4n+1$, by $p$.

Let $S$ denote the collection of all simply least residues for the modulus $a+bi=m$, with zero excluded, so that the count of the numbers contained in $S$ is $=p-1$. Let $x+yi$ denote any number of this system, and put $ax+by=\xi,ay-bx=\eta$. Then $\xi,\eta$ will be integers lying between the bounds $0$ and $p$, \emph{excluding the bounds}. For in the present case, where $a,b$ are relatively prime, the formulas of Art.~45, namely $\eta\equiv k\xi,\xi\equiv-k\eta\pmod p$, show that neither of the numbers $\xi,\eta$ can be $=0$ unless the other vanishes at the same time, so that $x=0,y=0$, a combination we have already excluded. Thus the criterion for a number $x+yi$ to be contained in $S$ is that the four numbers $\xi,\eta,p-\xi,p-\eta$ be positive.

We further observe that $\xi=\eta$ cannot hold for any such number. For it would follow that $p(x+y)=a(\xi+\eta)+b(\xi-\eta)=2a\xi$, which is absurd, since none of the factors $2,a,\xi$ is divisible by $p$. Similarly, the equation $p(x-y+a+b)=2a\xi+(a+b)(p-\xi-\eta)$ shows that $\xi+\eta=p$ cannot hold. Since the numbers $\xi-\eta,p-\xi-\eta$ must therefore be either positive or negative, we obtain from this a subdivision of the system $S$ into four collections $C,C',C'',C'''$, assigning the numbers as follows:
\[
\begin{array}{c|ll}
\text{Into collection}&\multicolumn{2}{c}{\text{Numbers for which}}\\\hline
C&\xi-\eta\text{ positive},&p-\xi-\eta\text{ positive}\\
C'&\xi-\eta\text{ positive},&p-\xi-\eta\text{ negative}\\
C''&\xi-\eta\text{ negative},&p-\xi-\eta\text{ negative}\\
C'''&\xi-\eta\text{ negative},&p-\xi-\eta\text{ positive}
\end{array}
\]

The criterion for a number of the collection $C$ is thus, strictly speaking, sixfold: the six numbers $\xi,\eta,p-\xi,p-\eta,\xi-\eta,p-\xi-\eta$ must be positive. But conditions $2,5$ and $6$ evidently already imply the others. The same applies to the collections $C',C'',C'''$, so that the complete criteria are threefold, namely:
\SourceLeaf{151}%
\[
\begin{array}{c|lll}
\text{For collection}&\multicolumn{3}{c}{\text{The following numbers must be positive}}\\\hline
C&\eta,&\xi-\eta,&p-\xi-\eta\\
C'&p-\xi,&\xi-\eta,&\xi+\eta-p\\
C''&p-\eta,&\eta-\xi,&\xi+\eta-p\\
C'''&\xi,&\eta-\xi,&p-\xi-\eta
\end{array}
\]

\Needspace{10\baselineskip}
Even without our pointing it out, anyone will readily understand that, in the geometric representation of complex numbers (see Art.~39), the numbers of the system $S$ lie inside the square whose sides join the points representing the numbers $0,a+bi,(1+i)(a+bi),i(a+bi)$, and that the subdivision of the system $S$ corresponds to the division of the square by its diagonals. Here, however, we have preferred to use purely arithmetical reasoning, leaving illustration by a geometric representation to the experienced reader for the sake of brevity.

\Needspace{22\baselineskip}
\Article{70}
If four complex numbers $r=x+yi,r'=x'+y'i,r''=x''+y''i,r'''=x'''+y'''i$ are related in such a way that $r'=m+ir,r''=m+ir'=(1+i)m-r,r'''=m+ir''=im-ir$, and the first, $r$, is supposed to belong to the collection $C$, then the others $r',r'',r'''$ will belong respectively to the collections $C',C'',C'''$. For, putting $\xi=ax+by,\eta=ay-bx,\xi'=ax'+by',\eta'=ay'-bx',\xi''=ax''+by'',\eta''=ay''-bx'',\xi'''=ax'''+by''',\eta'''=ay'''-bx'''$, we find
\[
\begin{aligned}
\eta&=p-\xi'=p-\eta''=\xi'''\\
\xi-\eta&=\xi'+\eta'-p=\eta''-\xi''=p-\xi'''-\eta'''\\
p-\xi-\eta&=\xi'-\eta'=\xi''+\eta''-p=\eta'''-\xi'''
\end{aligned}
\]
from which the truth of the theorem follows at once by the criteria. Since $r=m+ir'''$ holds again, it is easy to see that, if $r$ is supposed to belong to $C'$, the numbers $r',r'',r'''$ belong respectively to $C'',C''',C$; if it belongs to $C''$, they belong to $C''',C,C'$; and finally, if it belongs to $C'''$, they belong to $C,C',C''$.

At the same time, we conclude that the individual collections $C,C',C'',C'''$ contain equally many numbers, namely $\frac14(p-1)$.

\par\medskip\Needspace{4\baselineskip}\pdfbookmark[1]{Art. 37}{paper-five-native151-art37}\centerline{\textbf{37.}}\nobreak\smallskip
\textsc{Theorem.} \emph{Let $k$ denote an integer not divisible by $m$. If each number of the collection $C$ is multiplied by $k$, and the simply least residues of the products modulo
\SourceLeaf{152}%
$m$ are distributed among the collections $C,C',C'',C'''$, with the counts of the residues belonging to these individual collections denoted respectively by $c,c',c'',c'''$, then the character of the number $k$ with respect to the modulus $m$ will be $\equiv c'+2c''+3c'''\pmod4$.}

\Needspace{22\baselineskip}
\emph{Proof.} Let those $c$ least residues belonging to $C$ be $\alpha,\beta,\gamma,\delta$, etc.; next, let the $c'$ residues belonging to $C'$ be $m+i\alpha',m+i\beta',m+i\gamma',m+i\delta'$, etc.; further, let the $c''$ residues belonging to $C''$ be $(1+i)m-\alpha'',(1+i)m-\beta'',(1+i)m-\gamma'',(1+i)m-\delta''$, etc.; and finally, let the $c'''$ residues belonging to $C'''$ be $im-i\alpha''',im-i\beta''',im-i\gamma''',im-i\delta'''$, etc. Now consider four products, namely:
\begin{enumerate}
\renewcommand{\labelenumi}{\arabic{enumi})}
\item the product of all $\frac14(p-1)$ numbers constituting the collection $C$:
\item the product of the products obtained by multiplying each of these numbers by $k$;
\item the product of the least residues of these products, namely of the numbers $\alpha,\beta,\gamma,\delta$, etc., $m+i\alpha',m+i\beta'$, etc., etc.
\item the product of all $c+c'+c''+c'''$ numbers $\alpha,\beta,\gamma,\delta$, etc., $\alpha',\beta',\gamma',\delta'$, etc., $\alpha'',\beta'',\gamma'',\delta''$, etc.
\end{enumerate}
Denoting these four products in order by $P,P',P'',P'''$, we evidently have
\[
\begin{gathered}
P'=k^{\frac14(p-1)}P,\quad P'\equiv P'',\quad
P''\equiv P'''i^{c'+2c''+3c'''}\pmod m
\end{gathered}
\]
and therefore
\[
Pk^{\frac14(p-1)}\equiv P'''i^{c'+2c''+3c'''}\pmod m
\]

It is easy to see, however, that the numbers $\alpha',\allowbreak\beta',\allowbreak\gamma',\allowbreak\delta'$, etc., $\alpha'',\allowbreak\beta'',\allowbreak\gamma'',\allowbreak\delta''$, etc., $\alpha''',\allowbreak\beta''',\allowbreak\gamma''',\allowbreak\delta'''$, etc., all belong to the collection $C$, and are distinct both from one another and from the numbers $\alpha,\allowbreak\beta,\allowbreak\gamma,\allowbreak\delta$, etc., just as the latter themselves are distinct from one another. All these numbers taken together, disregarding their order, must therefore be exactly the same as all the numbers constituting the collection $C$. Hence we conclude that $P=P'''$, and consequently
\[
Pk^{\frac14(p-1)}\equiv Pi^{c'+2c''+3c'''}\pmod m
\]
Finally, since none of the individual factors of the product $P$ is divisible by $m$, we conclude that
\[
k^{\frac14(p-1)}\equiv i^{c'+2c''+3c'''}\pmod m
\]
so that $c'+2c''+3c'''$ will be the character of the number $k$ with respect to the modulus $m$. Q.~E.~D.

\SourceLeaf{153}%
\Article{72}
To apply the general theorem of the preceding article to the number $1+i$, we must subdivide the collection $C$ again into two smaller collections $G$ and $G'$. We shall assign to the collection $G$ those numbers $x+yi$ for which $ax+by=\xi$ is less than $\frac12p$, and to the other collection $G'$ those for which $\xi$ is greater than $\frac12p$. We shall denote the counts of the numbers contained in the collections $G,G'$ respectively by $g,g'$, so that $g+g'=\frac14(p-1)$.

The complete criterion for numbers belonging to $G$ is therefore that the three numbers $\eta,\xi-\eta,p-2\xi$ be positive. For the third condition for the collection $C$, requiring $p-\xi-\eta$ to be positive, is already implicitly contained in these conditions, since $p-\xi-\eta=(\xi-\eta)+(p-2\xi)$. Similarly, the complete criterion for numbers belonging to $G'$ consists in positive values of the three numbers $\eta,p-\xi-\eta,2\xi-p$.

It follows easily that the product of any number of the collection $G$ by the number $1+i$ belongs to the collection $C'''$. For, putting
\[
\begin{gathered}
(x+yi)(1+i)=x'+y'i,\quad\text{and}\quad ax'+by'=\xi',\ ay'-bx'=\eta',\ \text{we find}\\
\xi'=\xi-\eta,\quad\eta'-\xi'=2\eta,\quad p-\xi'-\eta'=p-2\xi
\end{gathered}
\]
that is, the criterion for the number $x+yi$ to belong to the collection $G$ is identical with the criterion for the number $x'+y'i$ to belong to the collection $C'''$.

In precisely the same way, it is shown that the product of any number of the collection $G'$ by $1+i$ belongs to the collection $C''$.

\Needspace{12\baselineskip}
Thus, if in the preceding article we give $k$ the value $1+i$, we have $c=0,c'=0,c''=g',c'''=g$, and the character of the number $1+i$ becomes $3g+2g'=\frac12(p-1)+g$. Since the characters of the numbers $i,-1$ are $\frac14(p-1),\frac12(p-1)$, the characters of the numbers $-1+i,-1-i,1-i$ will respectively be $\frac34(p-1)+g,g,\frac14(p-1)+g$. The whole matter now turns, therefore, on determining the number $g$.
